\section*{अध्याय \ref{ch:b2:amines} --- ऐमीन और नाइट्रोजन यौगिक}

\begin{solution}{exo:b2:amines:1}
प्राथमिक; द्वितीयक; तृतीयक; चतुष्क अमोनियम लवण; प्राथमिक (\omterm{def:b2:huckel:aromatic}{ऐरोमैटिक})।
\end{solution}

\begin{solution}{exo:b2:amines:2}
एथेनैमाइड $<$ ऐनिलीन $<$ अमोनिया $<$ मेथिलऐमीन (संयुग्मी अम्लों के $\mathrm pK_a$
$-0.6$, 4.60, 9.26, 10.66)।
\end{solution}

\begin{solution}{exo:b2:amines:3}
\textit{N}-मेथिलसाइक्लोहेक्सेनिमीन, \ce{C6H10=N-CH3}: दुर्बल अम्लीय माध्यम (pH 4--5) में ऐमीन और कीटोन,
बने जल को हटाते हुए (शुष्ककारक या \omterm{def:b2:liquid-vapour-diagrams:azeotrope}{स्थिरक्वाथी} आसवन)।
\end{solution}

\begin{solution}{exo:b2:amines:4}
ब्रोमोबेन्ज़ीन; बेन्ज़ोनाइट्राइल; फ़ीनॉल; 4-हाइड्रॉक्सीऐज़ोबेन्ज़ीन \ce{C6H5-N=N-C6H4-OH} (फ़ीनॉक्साइड के
\ce{O-} समूह के पैरा स्थान पर युग्मन)।
\end{solution}

\begin{solution}{exo:b2:amines:5}
$\mathrm pK_a = 9.80$: pH 7 पर अनुपात $[\ce{(CH3)3NH+}]/[\ce{(CH3)3N}] = 10^{2.8} \approx 630$:
प्रोटॉनित रूप प्रभावी है; नींबू का रस (pH 2 के निकट) व्यावहारिक रूप से इसे पूरा प्रोटॉनित कर देता है। किसी ऐमीन के
कार्बनिक विलयन को तनु जलीय अम्ल के साथ हिलाने से ऐमीन, अपने अमोनियम आयन के रूप में,
जल में चला जाता है; क्षारक उसे फिर मुक्त कर देता है।
\end{solution}

\begin{solution}{exo:b2:amines:6}
1-(साइक्लोहेक्स-1-ईन-1-इल)पिरोलिडीन। योग और जल की हानि के बाद नाइट्रोजन पर
धनात्मक आवेश होता है (एक इमीनियम आयन) और खोने के लिए कोई हाइड्रोजन नहीं; इसके बजाय सन्निकट
कार्बन से एक प्रोटॉन निकलता है, जिससे \ce{C=C} आबंध बनता है।
\end{solution}

\begin{solution}{exo:b2:amines:7}
प्रोपेन-2-ऐमीन के साथ बेन्ज़ैल्डिहाइड, या बेन्ज़िलऐमीन के साथ प्रोपेनोन, हर एक pH
लगभग 6 पर \ce{NaBH3CN} के साथ।
\end{solution}

\begin{solution}{exo:b2:amines:8}
प्रोपियोफ़ीनोन (1-फ़ेनिलप्रोपेन-1-ओन), \ce{C6H5COCH2CH3}: एक योग \omterm{def:b2:amines:imine}{इमीन} ऋणायन देता है,
जिसका जल-अपघटन कीटोन देता है।
\end{solution}

\begin{solution}{exo:b2:amines:9}
पोटैशियम थैलिमाइड + 1-ब्रोमोहेक्सेन: \textit{N}-हेक्सिलथैलिमाइड; फिर हाइड्रैज़ीन (या जल-अपघटन)
हेक्सेन-1-ऐमीन को मुक्त करता है। ऐल्किलित थैलिमाइड नाइट्रोजन के पास कोई उपलब्ध एकाकी युग्म नहीं होता (वह
दो कार्बोनिलों पर विस्थानीकृत है) और कोई \ce{N-H} नहीं बचता: उसका दूसरी बार ऐल्किलन नहीं हो सकता।
\end{solution}

\begin{solution}{exo:b2:amines:10}
ब्रोमीन जल ऐनिलीन का ऐमीनो समूह के ऑर्थो और पैरा स्थित तीनों स्थितियों पर ब्रोमीनन करता है:
2,4,6-ट्राइब्रोमोऐनिलीन। \omterm{def:b2:amines:diazonium}{डाइऐज़ोकरण}, फिर \ce{H3PO2}, ऐमीनो समूह को H से प्रतिस्थापित करता है: तीनों
ब्रोमीन अब एक-दूसरे के सापेक्ष 1,3,5 पर हैं।
\end{solution}

\begin{solution}{exo:b2:amines:11}\emergencystretch=3em
0--\qty{5}{\degreeCelsius} पर 4-नाइट्रोऐनिलीन + \ce{NaNO2} + \ce{HCl}: 4-नाइट्रोबेन्ज़ीनडाइऐज़ोनियम
क्लोराइड। दुर्बल क्षारीय विलयन में (फ़ीनॉक्साइड, प्रबल रूप से सक्रियित) फ़ीनॉल के साथ
\ce{O-} के पैरा स्थान पर युग्मन: 4-[(4-नाइट्रोफ़ेनिल)डाइऐज़ेनिल]फ़ीनॉल, \ce{O2N-C6H4-N=N-C6H4-OH}, एक नारंगी
रंजक।
\end{solution}

\begin{solution}{exo:b2:amines:12}
ट्राइमेथिलअमोनियम समूह एक क्षीण, भारी-भरकम निर्गामी समूह है और
\ce{CH3} वाली ओर के $\beta$ हाइड्रोजनों को अधिक अम्लीय और अधिक सुलभ बनाता है; क्षारक कम
प्रतिस्थापित कार्बन से हाइड्रोजन हटाता है: ब्यूट-1-ईन, कम प्रतिस्थापित ऐल्कीन (हॉफ़मान उत्पाद)।
\end{solution}

\begin{solution}{pb:b2:amines:1}\emergencystretch=3em
\textbf{1.} $\mathrm{{}^-O_3S{-}C_6H_4{-}NH_3^+}$: अम्लीय सल्फ़ोनिक समूह ने ऐमीन को प्रोटॉनित कर दिया है।
\textbf{2.} उभयाविष्ट आयन जल में अल्प-विलेय है; कार्बोनेट अमोनियम समूह को विप्रोटॉनित करता है,
जिससे विलेय सोडियम सल्फ़ैनिलेट बनता है।
\textbf{3.} \ce{NaNO2 + HCl -> HNO2 + NaCl}।
\textbf{4.}
\[
  \mathrm{{}^-O_3S{-}C_6H_4{-}NH_2 + HNO_2 + H^+ \to {}^-O_3S{-}C_6H_4{-}N_2^+ + 2\,H_2O} .
\]
\textbf{5.} गर्म होने पर डाइऐज़ोनियम लवण (फ़ीनॉल और \ce{N2} में) विघटित हो जाता है; नाइट्रस अम्ल
भी विघटित होता है।
\textbf{6.} $5.00/173.19 = \qty{28.9}{mmol}$: \ce{NaNO2} के \qty{28.9}{mmol}, अर्थात् \qty{1.99}{g}।
\textbf{7.} आधिक्य नाइट्रस अम्ल युग्मन भागीदार से अभिक्रिया कर लेगा; इसका पता स्टार्च--आयोडाइड
पत्र से लगाया जाता है और इसे यूरिया या सल्फ़ैमिक अम्ल से नष्ट किया जाता है।
\textbf{8.} इसका वलय डाइमेथिलऐमीनो समूह, एक शक्तिशाली दाता, द्वारा प्रबल रूप से सक्रियित है।
\textbf{9.} \ce{N(CH3)2} के पैरा स्थान पर: ऑर्थो/पैरा निर्देशन, और ऑर्थो स्थितियाँ
मेथिल समूहों द्वारा बाधित हैं।
\textbf{10.} $\mathrm{{}^-O_3S{-}C_6H_4{-}N{=}N{-}C_6H_4{-}N(CH_3)_2}$।
\textbf{11.} प्रबल अम्ल में डाइमेथिलऐमीनो समूह प्रोटॉनित हो जाता है: \ce{-NH(CH3)2+} वलय को निष्क्रिय करता है
और युग्मन रुक जाता है।
\textbf{12.} उत्पाद अपने प्रोटॉनित, आंशिक रूप से विलेय अम्ल रूप में बनता है; क्षारक में सोडियम
सल्फ़ोनेट, जो लवणीय विलयन में कम विलेय है, अवक्षेपित होता है।
\textbf{13.} इसका धनात्मक आवेश वलय पर विस्थानीकृत है और अंतस्थ नाइट्रोजन एक क्षीण
इलेक्ट्रॉनरागी केंद्र है; केवल \ce{-O-}, \ce{-OH} या \ce{-NR2} से सक्रियित वलय अभिक्रिया करते हैं।
\textbf{14.} क्षारक में पीला-नारंगी, अम्ल में लाल।
\textbf{15.} ऐज़ो समूह के एक नाइट्रोजन पर (प्रोटॉनित रूप डाइमेथिलऐमीनो समूह के साथ संयुग्मन से
स्थायीकृत है)।
\textbf{16.} लगभग $\mathrm pK_a \pm 1$: मोटे तौर पर 3.1 से 4.4 तक, $\mathrm pK_a = 3.47$ के आसपास।
\textbf{17.} दो वलय और ऐज़ो समूह एक लंबा संयुग्मित निकाय बनाते हैं: उच्चतम
भरे और निम्नतम खाली $\pi$ स्तरों के बीच का अंतराल छोटा है (\cref{prop:b2:huckel:gap}) और अवशोषण
दृश्य क्षेत्र में होता है।
\textbf{18.} प्रोटॉनन $\pi$ निकाय में आवेश का वितरण बदलता है और अंतराल को और संकरा
करता है: अवशोषण लंबी तरंगदैर्ध्यों की ओर खिसकता है, रंग पीले से लाल हो जाता है।
\textbf{19.} \qty{173.19}{g/mol}।
\textbf{20.} \qty{327.33}{g/mol}।
\textbf{21.} \qty{28.87}{mmol}।
\textbf{22.} $0.02887 \times 327.33 = \qty{9.45}{g}$।
\textbf{23.} डाइऐज़ोनियम लवण का विघटन, अपूर्ण युग्मन, मातृ द्रव और धोवन में
उत्पाद की विलेयता।
\textbf{24.} $0.80 \times 9.45 = \boldsymbol{\qty{7.56}{g}}$ मेथिल ऑरेंज।
\end{solution}
